\documentclass[12pt]{article}
\usepackage[utf8]{inputenc}
\usepackage[T1]{fontenc}
\usepackage{hyperref}
\title{Kgalagadi Language Resources for K-12 Curriculum Development}
\author{Compiled by Claude Sonnet 5.5 (claude-sonnet-5-5) for Andr\'as Kornai}
\date{}
\begin{document}
\maketitle

\section{Introduction}
Kgalagadi (SheKgalagadi; ISO 639-3 xkv; Glottocode kgal1244) is a Bantu language of Botswana, closely related to Tswana. It was only recently classified as a separate language rather than a Tswana dialect. Speaker estimates range from about 40,000 to 65,400. Wikitongues rates it as vigorous (EGIDS 6a). Search found no dedicated dictionary, standard orthography or school materials.

\section{Primary Educational Resources}
No primers, textbooks or school programmes for Kgalagadi were found in search results.

\section{Linguistic and Reference Materials}
\begin{itemize}
\item Wikipedia, ``Kgalagadi language'': overview. \url{https://en.wikipedia.org/wiki/Kgalagadi_language}
\item Omniglot, Kgalagadi (SheKgalagadi): alphabet and pronunciation information. According to the search summary, it also links a grammar sketch of the Shetjhauba variety. \url{https://www.omniglot.com/writing/kgalagadi.htm}
\item Endangered Languages Project record for Kgalagadi. It reports no standard orthography information and cites Andersson and Janson 1997, ``Sekgalagadi: Dialect or Language?'', and Hasselbring 2000, a sociolinguistic survey of Botswana's languages. \url{https://endangeredlanguages.com/node/86239}
\item Glosbe: crowd-sourced translation pages. They are not a vetted dictionary, so use with caution. \url{https://glosbe.com/xkv/lib}
\end{itemize}

\section{Religious and Community Resources}
None specific to Kgalagadi were found in search results.

\section{Digital and Audio Resources}
\begin{itemize}
\item Wikitongues, Kgalagadi (xkv): language page. \url{https://wikitongues.org/languages/xkv/}
\end{itemize}

\section{Curriculum Development Context}
No evidence of schooling in Kgalagadi was found. Botswana's school system uses Setswana and English. Because Kgalagadi is close to Tswana, Tswana materials could serve as a bridge. A curriculum would need an agreed orthography and community-developed materials. Government or academic orthography bodies were not checked.
\end{document}